\documentclass[12pt,a4paper]{article}

% -----------------------------------------------------------------------
% Packages
% -----------------------------------------------------------------------
\usepackage[T1]{fontenc}
\usepackage[utf8]{inputenc}
\usepackage[ngerman]{babel}
\usepackage{lmodern}
\usepackage{geometry}
\geometry{left=2.5cm, right=2.5cm, top=2.5cm, bottom=2.5cm}

\usepackage{booktabs}
\usepackage{longtable}
\usepackage{array}
\usepackage{tabularx}
\usepackage{enumitem}
\usepackage{titlesec}
\usepackage{parskip}
\usepackage{microtype}
\usepackage{hyperref}
\hypersetup{
    colorlinks=true,
    linkcolor=blue!70!black,
    urlcolor=blue!70!black,
    pdftitle={ICME4CRASH},
    pdfauthor={Konsortium ICME4CRASH},
}

% -----------------------------------------------------------------------
% Formatting helpers
% -----------------------------------------------------------------------
\newcolumntype{L}[1]{>{\raggedright\arraybackslash}p{#1}}

% -----------------------------------------------------------------------
\begin{document}

% =======================================================================
% Title
% =======================================================================
\begin{titlepage}
    \centering
    \vspace*{2cm}
    {\Huge\bfseries ICME4CRASH\par}
    \vspace{0.5cm}
    {\large\itshape Vollständig virtuelle, KI-gestützte Materialkarten für
    crashrelevante Multi-Material-Strukturen\par}
    \vspace{1.5cm}
    {\large Projektkonzept -- EU Horizon Europe RIA\par}
    \vfill
    {\normalsize Stand: April 2026}
\end{titlepage}

\tableofcontents
\clearpage

% =======================================================================
\section{Einleitung}
% =======================================================================

Die moderne Automobilindustrie setzt zunehmend auf eine Strategie des
Materialmix (\emph{Multi-Material-Mix}), um die strengen Anforderungen an
Fahrzeugsicherheit, Kraftstoffeffizienz und Gewichtsreduzierung zu erfüllen.
Dieser Ansatz kombiniert ultrahochfeste Stähle wie 22MnB5 mit
Hochleistungs-Aluminiumlegierungen (der Serien 6\textit{xxx} und
7\textit{xxx}). Die mechanische Integrität und das Crashverhalten dieser
Leichtbaukonstruktionen hängen jedoch stark von ihrer komplexen
mikrostrukturellen Entwicklung während der Fertigung ab~-- von
Abschreckprozessen bei Stahl bis hin zur Aushärtungskinetik bei Aluminium.

Traditionell erfordert die Charakterisierung des Versagensverhaltens für
Crash-Modelle wie \emph{LS-DYNA's GISSMO} oder \emph{ESI VPS's MMC}
umfangreiche experimentelle Tests für jeden spezifischen Material- und
Prozesszustand. Diese Tests sind teuer und zeitaufwendig und liefern in der
Regel nur die Materialdaten für den idealen Ausgangszustand. Insbesondere an
thermischen Verbindungen~-- die ohnehin aufgrund von Kerbwirkungen
hochkritisch sein können~-- können die Materialeigenschaften erheblich von
denen des Grundmaterials abweichen. Für diese geschwächten
Wärmeeinflusszonen liegen in der Regel keine Materialdatenblätter vor, was
bedeutet, dass die strukturelle Integrität in diesen Bereichen oft
überschätzt wird. Dies erfordert daher die Entwicklung rein virtueller
Materialkarten, die nicht nur die chemische Zusammensetzung, sondern auch
die Zeit-Temperatur-Zyklen berücksichtigen, denen das Material ausgesetzt
war.

Diese Forschungsarbeit schlägt einen robusten, vollständig virtuellen
Charakterisierungsrahmen vor, der auf \emph{Integrated Computational
Materials Engineering} (ICME) basiert. Durch die Nutzung thermokinetischer
Simulationen (z.\,B.\ \textbf{MatCalc}) zur Vorhersage von Phasen- und
Ausscheidungszuständen und deren Kopplung mit stochastischen \emph{Statistical
Volume Elements} (SVEs) wird die durch lokale mikrostrukturelle Variationen
bedingte inhärente Materialstreuung erfasst.

Unter Verwendung von Monte-Carlo-basierter Kristallplastizität
(CPFEM/CPFFT) leitet dieser Arbeitsablauf statistisch validierte
Versagensflächen über verschiedene Spannungszustände hinweg ab. Diese
Methode bietet einen skalierbaren \emph{Digital Twin} für das
Materialverhalten und ermöglicht die schnelle Erstellung hochpräziser
Materialkarten für heterogene Automobilstrukturen bei gleichzeitiger
erheblicher Reduzierung der Kosten für physikalische Prototypen.

% =======================================================================
\section{Forschungsansatz}
% =======================================================================

Der Forschungsansatz gliedert sich in eine durchgehende fünfstufige
ICME-Toolkette (\emph{Digital Thread} von der Legierung bis zur
Crash-Karte):

% -----------------------------------------------------------------------
\subsection{Stufe 1: Thermo-Kinetische Prozesssimulation -- Die
            \emph{Metallurgische DNA}}
% -----------------------------------------------------------------------

In diesem Schritt wird der Einfluss des Fertigungsprozesses (z.\,B.\ der
Temperaturzyklus beim Laserschweißen oder das Abkühlen im
Presshärtewerkzeug) auf die Chemie des Materials berechnet.

\textbf{Was simuliert wird:}
\begin{itemize}[leftmargin=*]
  \item Die zeitliche Entwicklung des Werkstoffgefüges.
  \item Welche Phasen entstehen (z.\,B.\ Austenit wandelt sich in Martensit
        oder Bainit um) und wie groß werden die Phasenanteile.
  \item Wie Ausscheidungen (Präzipitate) in einer 6\textit{xxx}-Aluminiumlegierung
        während der Warmauslagerung oder durch die Schweißhitze wachsen.
\end{itemize}

\textbf{Methodik:} CALPHAD (\emph{Calculation of Phase Diagrams}) gekoppelt
mit kinetischen Ratengleichungen. Die Methodik löst thermodynamische
Gleichgewichte (Energieminimierung) und berechnet Diffusionsgeschwindigkeiten.

\textbf{Software-Optionen:}
\begin{itemize}[leftmargin=*]
  \item \textbf{MatCalc} -- Der Goldstandard aus Österreich/Wien; extrem stark
        bei Ausscheidungskinetik, perfekter Fit für das MCL Leoben.
  \item Thermo-Calc (mit Zusatzmodul TC-PRISMA).
  \item JMatPro -- Sehr industrienah; liefert direkt mechanische Basiswerte.
\end{itemize}

\textbf{Output für Stufe 2:} Phasenanteile (z.\,B.\ 80\,\% Martensit,
20\,\% Bainit), Versetzungsdichte und Teilchengrößen.

% -----------------------------------------------------------------------
\subsection{Stufe 2: Gefüge-Generierung -- Der stochastische digitale
            Zwilling}
% -----------------------------------------------------------------------

Die abstrakten Zahlenwerte aus Stufe~1 (Prozentsätze und Radien) werden in
ein echtes, geometrisches 3D-Modell übersetzt -- das sogenannte
\emph{SVE (Statistical Volume Element)}.

\textbf{Was simuliert wird:}
\begin{itemize}[leftmargin=*]
  \item Ein virtueller, dreidimensionaler Würfel des Mikrokosmos wird
        gebaut und mit Körnern gefüllt, deren Größe, Form und
        kristallographische Orientierung (Textur) der Realität entsprechen.
  \item Die in Stufe~1 berechneten Phasen werden in dieses Gitter eingefüllt.
  \item Um die Streuung der Natur abzubilden, werden oft Dutzende leicht
        unterschiedlicher SVE-Würfel für denselben Materialzustand generiert.
\end{itemize}

\textbf{Methodik:} Voronoi-Tesselation oder Monte-Carlo-Potts-Modelle.

\textbf{Software-Optionen:}
\begin{itemize}[leftmargin=*]
  \item \textbf{Neper} -- Open Source; extrem mächtig für komplexe
        Voronoi-Strukturen.
  \item \textbf{DREAM.3D} -- Open Source; Standard-Tool zur Rekonstruktion
        von Gefügen aus EBSD-Messdaten.
  \item GeoDict / Digimat-FE -- Kommerzielle Lösungen für Faserverbunde oder
        komplexe Defekte (Poren).
\end{itemize}

\textbf{Output für Stufe 3:} Ein 3D-Geometrie-Netz (Mesh oder Voxel-Grid)
des Gefüges, in dem jedem Element eine Phase und eine Kristallorientierung
zugewiesen ist.

% -----------------------------------------------------------------------
\subsection{Stufe 3: Mikromechanische Prüfung -- Das virtuelle Labor}
% -----------------------------------------------------------------------

Das Herzstück der ICME-Kette: der SVE-Würfel aus Stufe~2 wird virtuell in
alle Richtungen gezogen, gedrückt und geschert.

\textbf{Was simuliert wird:}
\begin{itemize}[leftmargin=*]
  \item Die makroskopische Last wird auf das Kristallgitter übertragen.
  \item Bewegung von Versetzungen (Gitterbaufehler) durch das Material.
  \item Lokale Spannungsstauungen an Korngrenzen; plastisches Fließen weicher
        Körner (Ferrit) gegenüber harten Körnern (Martensit).
  \item Bestimmung der Dehnung, bei der das Gitter lokal reißt
        (Schädigungsinitiation durch Porenwachstum oder Scherbänder).
\end{itemize}

\textbf{Methodik:} CPFFT (\emph{Crystal Plasticity Fast Fourier Transform})
oder CPFEM (\emph{Crystal Plasticity Finite Element Method}).

\textbf{Software-Optionen:}
\begin{itemize}[leftmargin=*]
  \item \textbf{DAMASK} (\emph{Düsseldorf Advanced Material Simulation Kit})
        -- entwickelt vom MPIE; das Nonplusultra der Open-Source-Kristallplastizität.
  \item Abaqus mit speziellen User-Material-Subroutinen (UMAT).
  \item Simcenter~3D / MultiMechanics -- kommerzieller Ansatz von Siemens.
\end{itemize}

\textbf{Output für Stufe 4:} Hunderte virtuelle Spannungs-Dehnungs-Kurven
für verschiedenste Triaxialitäten (Zug, Druck, Scherung) inklusive einer
Datenwolke lokaler Versagenszeitpunkte (Streuung).

% -----------------------------------------------------------------------
\subsection{Stufe 4: Homogenisierung \& KI-Surrogate Modeling -- Der
            Übersetzer}
% -----------------------------------------------------------------------

Die Mikrodaten aus Stufe~3 werden verdichtet, da ein Makro-Crash-Solver
(LS-DYNA) direkt damit nicht umgehen kann.

\textbf{Was hier geschieht:} Data Science statt physikalischer Simulation.
Ein Surrogate Model wird trainiert, das für jede beliebige Makro-Dehnung und
Temperatur sofort die zugehörige Makro-Spannung liefert, ohne die
Kristallplastizität erneut berechnen zu müssen.

\textbf{Methodik:} Volumenmittelung (\emph{Volume Averaging}) der
Spannungs-Dehnungs-Felder; Einsatz von Machine Learning (Neuronale Netze)
oder Response Surface Methods zur Erstellung einer durchgehenden
Versagensfläche (\emph{Failure Surface}).

\textbf{Software-Optionen:}
\begin{itemize}[leftmargin=*]
  \item Python-Ökosystem: PyTorch, TensorFlow, Scikit-learn.
  \item Altair HyperStudy / LS-OPT -- für automatisierte Sensitivitätsanalysen.
\end{itemize}

\textbf{Output für Stufe 5:} Mathematisch saubere Kurvenscharen
(deterministisch oder mit Konfidenzintervallen), frei von
mikroskopischem Rauschen.

% -----------------------------------------------------------------------
\subsection{Stufe 5: Materialkarten-Kalibrierung -- Crash Readiness}
% -----------------------------------------------------------------------

Die homogenisierten, virtuellen Kurven werden in die von OEMs geforderten
Materialkarten überführt.

\textbf{Was simuliert wird:} Parameter-Identifikation (\emph{Reverse
Engineering}). Ein Algorithmus variiert die Parameter der GISSMO- oder
MMC-Gleichungen, bis die analytische Karte exakt das Verhalten der
KI-Daten aus Stufe~4 widerspiegelt. Besonders wichtig ist die
\textbf{Elementgrößen-Regularisierung}: Der Algorithmus berechnet, wie die
Karte reagieren muss, wenn das Finite-Element in der Crashsimulation 1\,mm
groß ist im Vergleich zu 5\,mm.

\textbf{Methodik:} Nicht-lineare Optimierungsalgorithmen (z.\,B.\ Genetische
Algorithmen, Levenberg-Marquardt) zur Minimierung der Fehlerquadratsumme
zwischen virtueller Target-Kurve und Modell-Kurve.

\textbf{Software-Optionen:}
\begin{itemize}[leftmargin=*]
  \item \textbf{LS-OPT} -- Der Standard für LS-DYNA-Karten.
  \item Visual-Environment / PAM-OPT -- für ESI PAM-CRASH.
  \item MATLAB / Python mit Optimierungs-Toolboxen für maximale Kontrolle.
\end{itemize}

% =======================================================================
\section{Projektziele}
% =======================================================================

Für einen RIA-Antrag dieser Größenordnung müssen die \emph{Objectives} eine
perfekte Balance zwischen wissenschaftlicher Exzellenz und industriellem
Impact finden. In der EU-Nomenklatur werden sie unterteilt in
\emph{Scientific Objectives} (SO), \emph{Technological Objectives} (TO) und
\emph{Impact Objectives} (IO).

% -----------------------------------------------------------------------
\subsection{Wissenschaftliche Ziele (Scientific Objectives -- SO)}
% -----------------------------------------------------------------------

\begin{description}[leftmargin=*, style=nextline]
  \item[SO\,1: Physikbasierte Versagensvorhersage auf der Mesoskala]
    Entwicklung eines stochastischen Multiskalen-Frameworks, das die
    kristallplastische Verformung und Schädigungsinitiation unter
    Berücksichtigung von prozessinduzierten Defekten (z.\,B.\ Mikroporosität,
    lokale Entfestigung in der HAZ) mit einer Genauigkeitssteigerung von
    25\,\% gegenüber rein empirischen Modellen vorhersagt.

  \item[SO\,2: Skalenübergreifende Schadensmechanik]
    Erforschung der Korrelation zwischen mikroskopischen
    Versagensmechanismen (\emph{Void Nucleation/Growth}) und makroskopischen
    Versagensindikatoren (Triaxialität/Lode-Winkel) für mindestens drei neue
    Leichtbau-Materialkombinationen (AHSS-Stahl, 6\textit{xxx}/7\textit{xxx}
    Aluminium, Mischbau).

  \item[SO\,3: Physics-Informed AI (PINN)]
    Entwicklung neuartiger, physik-informierter neuronaler Netze zur
    Homogenisierung, die sicherstellen, dass die KI-generierten Materialkarten
    fundamentale thermodynamische und mechanische Erhaltungssätze einhalten.
\end{description}

% -----------------------------------------------------------------------
\subsection{Technologische Ziele (Technological Objectives -- TO)}
% -----------------------------------------------------------------------

\begin{description}[leftmargin=*, style=nextline]
  \item[TO\,1: Automatisierte ICME-Pipeline]
    Implementierung einer lückenlosen, digitalen Toolkette, die den
    Datenfluss von der thermo-kinetischen Simulation (MatCalc) über die
    SVE-Generierung (MCL) bis zur fertigen Crash-Karte (ViF) automatisiert
    und die Kalibrierungszeit um \textbf{über 80\,\%} reduziert (von Wochen
    auf Tage).

  \item[TO\,2: Interoperabilität \& Standardisierung]
    Etablierung einer auf dem \textbf{VMAP-Standard} basierenden
    Schnittstellen-Architektur, die den nahtlosen Datentransfer zwischen den
    unterschiedlichen Software-Silos der Partner (MatCalc, DAMASK, LS-DYNA,
    PAM-CRASH) ermöglicht.

  \item[TO\,3: Virtuelle Validierungsumgebung]
    Schaffung eines \emph{Virtual Material Labors}, das in der Lage ist, die
    für eine GISSMO-Karte notwendigen Kurvenscharen rein virtuell zu erzeugen,
    wobei lediglich 3--5 physische "`Ankerversuche"' zur finalen Absicherung
    benötigt werden.
\end{description}

% -----------------------------------------------------------------------
\subsection{Impact-Ziele (Impact Objectives -- IO)}
% -----------------------------------------------------------------------

\begin{description}[leftmargin=*, style=nextline]
  \item[IO\,1: Reduktion physischer Ressourcen]
    Demonstration einer Einsparung von \textbf{50--70\,\% der physischen
    Materialcharakterisierungstests} bei der Einführung neuer
    Werkstoff-Prozess-Kombinationen in die Fahrzeugserienentwicklung.

  \item[IO\,2: Beschleunigung der Time-to-Market]
    Verkürzung der Entwicklungszyklen für sicherheitskritische
    Karosseriestrukturen um mindestens 20\,\% durch die Vorverlegung der
    Materialabsicherung in die frühe digitale Phase (Soll-Ist-Abgleich am
    Demonstrator TRL~5).

  \item[IO\,3: Stärkung der europäischen Souveränität]
    Etablierung eines europaweiten Standards für die virtuelle
    Materialcharakterisierung, der die Abhängigkeit von proprietären,
    außereuropäischen Datenbanken verringert und die Wettbewerbsfähigkeit von
    Playern wie Gestamp, Magna und Stellantis sichert.
\end{description}

% -----------------------------------------------------------------------
\subsection{Executive Summary -- Kernziel}
% -----------------------------------------------------------------------

\begin{quote}
  \emph{"`Das Projekt \textbf{[PROJEKTNAME]} wird die erste europäische,
  KI-gestützte ICME-Plattform etablieren, die den `Digital Thread' von der
  metallurgischen Chemie bis zur Fahrzeug-Sicherheit schließt. Durch die
  Integration prozessbedingter Defekte in die virtuelle Charakterisierung
  wird die Lücke zwischen Labor und Realität geschlossen, wodurch physische
  Crash-Validierungsschleifen drastisch reduziert und die Einführung
  innovativer Leichtbauwerkstoffe radikal beschleunigt werden."'}
\end{quote}

\textbf{Tipp -- \emph{The Power of 5 OEMs}:} Nutzen Sie die Beteiligung von
5~OEMs als \emph{Cross-Validation Objective}: Nachweis der universellen
Anwendbarkeit der entwickelten Methodik durch die parallele Validierung an
5~unterschiedlichen fahrzeugspezifischen Use-Cases (von der Batteriewanne
bei Audi bis zum Fahrwerksteil bei Ford Otosan).

% =======================================================================
\section{Mögliches Konsortium und Rollenverteilung}
% =======================================================================

\begin{longtable}{L{3cm} L{4cm} L{7.5cm}}
  \toprule
  \textbf{Sektor} & \textbf{Partner} & \textbf{Spezifische Rolle im
  RIA-Projekt} \\
  \midrule
  \endhead
  \bottomrule
  \endfoot

  Lead \& AI
    & Virtual Vehicle (AT)
    & Gesamtkoordination, Entwicklung der KI-Surrogate-Modelle (Stufe~4)
      und Systemintegration. \\[4pt]

  Software \& Kinetik
    & \textbf{MatCalc} (AT)
    & Eigener Software-Partner: Berechnung der Phasenwandlung und
      Ausscheidungskinetik (Stufe~1). \\[4pt]

  Physics \& ICME
    & \textbf{MCL Leoben} (AT)
    & Leitung der Multiskalen-Kette, SVE-Generierung und Kristallplastizität
      (Stufe~2 \& 3). \\[4pt]

  Micromechanics
    & \textbf{UNN} (UK)
    & Spezialisierte Schadensmodellierung für Aluminium-Leichtbau und neue
      CP-Modelle (Stufe~3). \\[4pt]

  Joining Tech
    & \textbf{LWF Paderborn} (DE)
    & Experte für thermische Behandlung und metallographische Präparation
      (Der "`Datenlieferant der Wahrheit"'). \\[4pt]

  Process \& Mfg.
    & \textbf{CTAG} (ES)
    & Überwachung der Fertigungsprozesse; Kopplung von Prozess-Sensorik mit
      Simulations-Input. \\[4pt]

  Large-Scale Test
    & \textbf{CIDAUT} (ES)
    & Dynamische Bauteiltests (Crash) zur finalen Validierung der virtuellen
      Karten (TRL~5). \\[4pt]

  Material
    & \textbf{voestalpine} (AT) \& \textbf{ArcelorMittal} (LU/FR),
      \textbf{AMAG} (AT)
    & Bereitstellung der "`metallurgischen DNA"' (Chemie, Gefüge-Basics) für
      neue Stahlgüten (z.\,B.\ PHS). \\[4pt]

  Tier-1 Supplier
    & \textbf{Gestamp} (ES) \& \textbf{Magna} (AT)
    & Use-Case-Provider für Strukturbauteile (B-Säule, Längsträger).
      Brücke zur Produktion. \\[4pt]

  OEM (End User)
    & \textbf{VW / Audi / Volvo / Ford Otosan / CRF} (Stellantis)
    & Use-Case-Provider: Definition der Crash-Lastfälle, Validierung am
      Gesamtfahrzeugmodell, Abnahme der Workflows. \\
\end{longtable}

% =======================================================================
\section{Workpackage-Struktur}
% =======================================================================

Das Projekt umfasst ein Budget von \textbf{6--8~Millionen Euro} mit ca.
14--16~Partnern. Die WP-Struktur bildet den Digital Thread von der Legierung
bis zum Gesamtfahrzeug-Crash ab.

% -----------------------------------------------------------------------
\subsection{WP\,1: Projektmanagement \& Strategische Koordination}
% -----------------------------------------------------------------------
\textbf{Lead:} Virtual Vehicle (ViF)

Administrative Leitung, Finanzcontrolling, IP-Management und Koordination
zwischen den 15+ Partnern. Aufbau des \emph{Project Executive Boards} mit
Vertretern der OEMs (Audi, Volvo, Ford, etc.).

% -----------------------------------------------------------------------
\subsection{WP\,2: Werkstoff-Engineering \& Prozess-Definition}
% -----------------------------------------------------------------------
\textbf{Lead:} voestalpine / ArcelorMittal (Stahl) \& Aluminium-Lieferant
(z.\,B.\ AMAG/Hydro)

\textbf{Aufgaben:}
\begin{itemize}[leftmargin=*]
  \item Definition der Legierungsvarianten (High-Strength Steel \& Cast
        Aluminum).
  \item Festlegung der industriellen Randbedingungen (Wärmebehandlungszyklen
        T4/T6, Presshärte-Parameter).
  \item Bereitstellung von Master-Chargen für das gesamte Konsortium.
\end{itemize}

% -----------------------------------------------------------------------
\subsection{WP\,3: Experimentelle Hochdurchsatz-Charakterisierung}
% -----------------------------------------------------------------------
\textbf{Lead:} LWF Paderborn

\textbf{Aufgaben:}
\begin{itemize}[leftmargin=*]
  \item \emph{Thermisches Processing:} Durchführung der
        Wärmebehandlungszyklen (Guss-Abkühlung, HAZ-Simulation).
  \item \emph{Metallographie-Zentrum:} Erzeugung von Schliffbildern zur
        Ermittlung von DAS (Dendritenarmabstand), Phasenanteilen und
        Porosität.
  \item \emph{Härte-Mapping:} Korrelation der Mikrostruktur mit mechanischen
        Basiseigenschaften.
\end{itemize}

% -----------------------------------------------------------------------
\subsection{WP\,4: Multiskalen-Physik \& Virtuelles Labor}
% -----------------------------------------------------------------------
\textbf{Lead:} MCL Leoben (Wissenschaft) \& MatCalc (Software)

\textbf{Aufgaben:}
\begin{itemize}[leftmargin=*]
  \item \emph{Stufe~1 (Kinetics):} MatCalc-Simulation der Phasenbildung und
        Ausscheidungskinetik.
  \item \emph{Stufe~2 (SVE):} MCL baut digitale 3D-Gefüge basierend auf den
        Schliffbildern aus WP\,3.
  \item \emph{Stufe~3 (CP):} Kristallplastizität (MCL/UNN) zur Vorhersage
        des lokalen Versagens (Kristall-Bruchmechanik).
  \item Partner: UNN -- Fokus auf Aluminium-Schädigung.
\end{itemize}

% -----------------------------------------------------------------------
\subsection{WP\,5: KI-Homogenisierung \& Tool-Integration}
% -----------------------------------------------------------------------
\textbf{Lead:} Virtual Vehicle (ViF)

\textbf{Aufgaben:}
\begin{itemize}[leftmargin=*]
  \item \emph{Surrogate Modeling:} Training von Neuronalen Netzen mit den
        Daten aus WP\,4.
  \item \emph{Mapping-Algorithmen:} Transfer von Guss-Parametern (Porosität)
        auf das Crash-Mesh.
  \item \emph{Software-Schnittstellen:} Integration in kommerzielle Solver
        (LS-DYNA / PAM-CRASH) mittels VMAP-Standard.
\end{itemize}

% -----------------------------------------------------------------------
\subsection{WP\,6: Demonstratoren \& System-Validierung (TRL\,5/6)}
% -----------------------------------------------------------------------
\textbf{Lead:} CIDAUT / CTAG (Experiment) \& Gestamp / Magna (Fertigung)

\textbf{Aufgaben:}
\begin{itemize}[leftmargin=*]
  \item Herstellung realer Bauteil-Demonstratoren (z.\,B.\ Gigacasting-Knoten
        oder B-Säule).
  \item \emph{Crash-Validierung:} Durchführung von Fallturm- und
        Schlittentests.
  \item Vergleich der Vorhersagegenauigkeit: Virtuelle ICME-Karte vs.
        Realer Crash.
\end{itemize}

% -----------------------------------------------------------------------
\subsection{WP\,7: Industrielle Use Cases \& OEM-Akzeptanz}
% -----------------------------------------------------------------------
\textbf{Lead:} OEMs (VW/Audi, Volvo, Ford Otosan, CRF)

\textbf{Aufgaben je Partner:}
\begin{itemize}[leftmargin=*]
  \item \textbf{Audi/VW:} Fokus auf Großserien-Stahlstrukturen \&
        Batterie-Crash.
  \item \textbf{Volvo:} Fokus auf Gigacasting-Strukturen und Virtual Testing
        Standards.
  \item \textbf{Ford Otosan:} Fokus auf schwere Nutzfahrzeugstrukturen
        (andere Dicken/Lasten).
  \item \textbf{CRF:} Multi-Material-Mix in der Kleinwagen-Plattform.
  \item \textbf{Magna/Gestamp:} Bewertung der prozessseitigen Robustheit.
\end{itemize}

% -----------------------------------------------------------------------
\subsection{WP\,8: Exploitation, Standardisierung \& Dissemination}
% -----------------------------------------------------------------------
\textbf{Lead:} ESI Group (oder ViF)

Kommerzielle Verwertung (Software-Plugins), wissenschaftliche Publikationen
und Training der nächsten Generation von Ingenieuren (\emph{Young Researcher
Programme}).

% =======================================================================
\section{Logischer Datenfluss (Gantt-Logik)}
% =======================================================================

\begin{description}[leftmargin=*]
  \item[Monat 1--12:] WP\,2 \& 3 liefern die physischen und
        metallographischen Daten (Schliffbilder).
  \item[Monat 6--24:] WP\,4 nutzt diese Daten für die tiefen physikalischen
        Modelle (MCL/MatCalc).
  \item[Monat 18--36:] WP\,5 trainiert die KI und baut die Schnittstellen.
  \item[Monat 24--42:] WP\,6 \& 7 validieren die Ergebnisse an echten
        Bauteilen; OEMs implementieren die Workflows.
\end{description}

% =======================================================================
\section{Budget-Allokation (Beispiel bei 8\,Mio.\ €)}
% =======================================================================

\begin{tabular}{L{6cm} r l}
  \toprule
  \textbf{Workpackage} & \textbf{Anteil} & \textbf{Betrag} \\
  \midrule
  WP\,1 (Management)                   & 5\,\%  & \~{}400\,k€ \\
  WP\,2 \& 3 (Material/Experiment)     & 20\,\% & \~{}1.600\,k€ \\
  WP\,4 (Physik/ICME)                  & 30\,\% & \~{}2.400\,k€ \\
  WP\,5 (KI/Software)                  & 20\,\% & \~{}1.600\,k€ \\
  WP\,6 \& 7 (Validation/OEMs)         & 20\,\% & \~{}1.600\,k€ \\
  WP\,8 (Exploitation)                 & 5\,\%  & \~{}400\,k€ \\
  \bottomrule
\end{tabular}

% =======================================================================
\section{Wichtige Meilensteine (Milestones)}
% =======================================================================

\begin{description}[leftmargin=*]
  \item[MS\,1 (M12):] Master-Schliffbild-Datenbank für Al-Guss \& PHS-Stahl
        verfügbar.
  \item[MS\,2 (M24):] Erste KI-basierte Materialkarte erfolgreich vom SVE
        zum Makro-Modell konvergiert.
  \item[MS\,3 (M36):] Demonstrator-Crashversuch bestätigt Fehlerreduktion
        von $>30\,\%$ gegenüber Standardmodellen.
\end{description}

% =======================================================================
\section{V-MAT-2030 -- Projektpitch}
% =======================================================================

\subsection*{The Digital Material Factory -- Unlocking the Future of
  Automotive Lightweighting through AI-Driven ICME}

\textbf{Die Vision.}
Europa steht vor der Revolution der \emph{Gigacastings} und des hochfesten
Mischbaus. Doch die traditionelle Materialcharakterisierung ist am Ende: Sie
ist zu langsam, zu teuer und ignoriert die "`schmutzige Realität"' der
Fertigung (Porosität, lokale Abkühlraten). \textbf{V-MAT-2030} schließt
diesen Digital Thread von der Legierung bis zum Crash-Test~-- rein virtuell,
physikalisch fundiert und KI-beschleunigt.

\textbf{Der Workflow: Vom Atom zum Euro-NCAP.}
Wir ersetzen hunderte physische Tests durch eine automatisierte
5-Stufen-Toolkette:
\begin{enumerate}[leftmargin=*]
  \item \emph{Metallurgische DNA (MatCalc):} Simulation der Phasenbildung
        basierend auf der Thermodynamik.
  \item \emph{Digitaler Zwilling des Gefüges (LWF \& MCL):} Verknüpfung
        realer Schliffbilder (DAS, Porosität) mit statistischen
        3D-Volumenelementen (SVE).
  \item \emph{Virtuelles Labor (MCL \& UNN):} Kristallplastizität simuliert
        das Versagen dort, wo keine Probe hinkommt.
  \item \emph{KI-Translation (ViF):} Machine Learning komprimiert Terabytes
        an Mikro-Daten in performante Makro-Materialkarten.
  \item \emph{Industrielle Validierung (CTAG, CIDAUT, OEMs):} Absicherung an
        realen Megacastings und B-Säulen.
\end{enumerate}

\textbf{Der wissenschaftliche Fokus: Mikrostruktur-DNA.}
Der Kern des Projekts ist die Erkenntnis, dass das Gefüge die Mechanik
diktiert. Wir nutzen hochauflösende Metallographie, um den
Dendritenarmabstand (DAS) und die Phasenmorphologie als primäre
Input-Geber für die Simulation zu etablieren.

\textbf{Das \emph{Champions League} Konsortium.}
Ein 8-Millionen-Euro-Verbund, der 40\,\% des europäischen Marktes
repräsentiert:
\begin{itemize}[leftmargin=*]
  \item RTOs \& Academy: Virtual Vehicle (Lead), MCL Leoben, LWF Paderborn,
        UNN, CTAG, CIDAUT.
  \item Software-Enabler: MatCalc, ESI Group.
  \item Material-Provider: voestalpine, ArcelorMittal, AMAG/Hydro.
  \item Use-Case-Provider \& OEMs: VW/Audi, Volvo, Ford Otosan, CRF
        (Stellantis), Magna, Gestamp.
\end{itemize}

\textbf{Strategische Ziele \& Impact (KPIs):}
\begin{itemize}[leftmargin=*]
  \item \textbf{Time-to-Market:} Reduktion der Materialabsicherung um
        $-80\,\%$.
  \item \textbf{Kosten:} Einsparung durch Wegfall von ca.\ $90\,\%$ der
        physischen Charakterisierungstests.
  \item \textbf{Leichtbau:} Ermöglicht $-15\,\%$ Gewichtseinsparung bei
        Megacastings durch präzise Vorhersage lokaler Eigenschaften statt
        hoher Sicherheitsfaktoren.
  \item \textbf{Nachhaltigkeit:} Erstmalige sichere Integration von
        Sekundär-Aluminium in crashrelevante Strukturen durch virtuelle
        Fehler-Modellierung.
\end{itemize}

% =======================================================================
\section{Strategische Positionierungsoptionen}
% =======================================================================

\subsection{Option~1: Der \emph{Physics-First}-Weg (Klassische RIA)}

\begin{description}[leftmargin=*, style=nextline]
  \item[Fokus:] Tiefgreifende wissenschaftliche Exzellenz; Lösung der
        Bruchmechanik auf Kornebene.
  \item[Kernargument:] "`Wir verstehen die Physik so gut, dass wir keine
        Experimente mehr brauchen."'
  \item[Stärke:] Extrem hohe Bewertung im Bereich "`Excellence"'. Gut für
        UNN und MCL.
  \item[Passt zu Calls, die:] Radikale Innovationen und \emph{Breakthrough
        Science} fordern.
\end{description}

\subsection{Option~2: Der \emph{Pragmatic-AI}-Weg (Digitaler Schatten)}

\begin{description}[leftmargin=*, style=nextline]
  \item[Fokus:] Datengetriebene Korrelation von Schliffbildern zu
        Materialkarten.
  \item[Kernargument:] "`Wir nutzen die KI, um das Expertenwissen der
        Metallographie (LWF) in Sekunden auf das ganze Auto zu übertragen."'
  \item[Stärke:] Sehr hohe Bewertung bei "`Impact"' und "`Exploitation"'. Gut
        für ViF, MatCalc und die Softwarehäuser (ESI).
  \item[Passt zu Calls, die:] Digitalisierung der Industrie, KI in der
        Produktion oder "`Twin Transition"' zum Thema haben.
\end{description}

\subsection{Option~3: Der \emph{Gigacasting-Sovereignty}-Weg
           (Application-Driven)}

\begin{description}[leftmargin=*, style=nextline]
  \item[Fokus:] Die größte aktuelle Herausforderung der europäischen
        Automotive-Industrie: Megacastings.
  \item[Kernargument:] "`Ohne unsere Toolkette kann Europa den Vorsprung von
        Tesla bei großen Gussstrukturen nicht einholen, ohne die Sicherheit
        zu gefährden."'
  \item[Stärke:] Höchste politische Relevanz; zieht die OEMs (Volvo, Audi,
        Ford) massiv an.
  \item[Passt zu Calls, die:] Mobilität, Zero-Emission-Fahrzeuge oder die
        Stärkung der europäischen Lieferkette (IAM4EU) adressieren.
\end{description}

\end{document}
